\chapter{Evaluate monadic programs}
\source{Sample Q4, p. 4; Monads lecture, do notation and List Monad.}
Use this supplied function for Questions 4.1--4.2. Do not reimplement it here.
\par\noindent\begin{minipage}{\linewidth}
\begin{lstlisting}
bundle :: Monad m => [m a] -> m [a]
bundle [] = return []
bundle (u:us) = do
  x <- u
  xs <- bundle us
  return (x:xs)
\end{lstlisting}
\end{minipage}\par

\task{4.1. Values and types}
Evaluate each expression and give its whole type. Leave function values such as
\code{id} unevaluated when appropriate. All parts are independent.
\begin{enumerate}[label=(\alph*)]
\item \code{bundle [Just 'a', Just 'c']}
\answer{[Result and type.]}
\item \code{bundle [Just not, Nothing]}
\answer{[Result and type.]}
\item \code{bundle [[] :: [Bool]]}
\answer{[Result and type.]}
\item \code{bundle ([] :: [[Bool]])}
\answer{[Result and type; distinguish this from part (c).]}
\item \code{bundle [[True,False], [True]]}
\answer{[All branches in order; whole type.]}
\item \code{bundle [[id,not], [const False]]}
\answer{[Function lists in order; then apply each returned function to True.]}
\end{enumerate}

\clearpage
\task{4.2. Translate and trace}
Rewrite the following without do notation, using binds and lambdas. Give each
bound variable's type and the final result. Explain the effect of the middle line.
\par\noindent\begin{minipage}{\linewidth}
\begin{lstlisting}
do x <- Just (6 :: Int)
   return (x+1)
   return (x*3)
\end{lstlisting}
\end{minipage}\par
\par\noindent\begin{minipage}{\linewidth}
\begin{codeanswer}
-- YOUR TRANSLATION
\end{codeanswer}
\end{minipage}\par
\answer{[Types, result, explanation.]}

\task{4.3. Implement safe monadic computation}
\code{safeQuot x y} returns \code{Nothing} for divisor zero and otherwise
\code{Just (x `div` y)}. \code{divideChain} divides the starting value by each
divisor left-to-right, failing on any zero; an empty divisor list succeeds with
the starting value. Use explicit recursion for \code{divideChain}.
\par\noindent\begin{minipage}{\linewidth}
\begin{codeanswer}
safeQuot :: Int -> Int -> Maybe Int
safeQuot = undefined  -- YOUR CODE

divideChain :: Int -> [Int] -> Maybe Int
divideChain = undefined  -- YOUR CODE
\end{codeanswer}
\end{minipage}\par

\task{4.4. List choices and a generic fold}
Implement \code{allPairs} using list do notation, with left choices outermost.
Implement \code{foldActions} as a monadic left fold using recursion and binds;
do not use library folds, \code{mapM}, or \code{sequence}.
\par\noindent\begin{minipage}{\linewidth}
\begin{codeanswer}
allPairs :: [a] -> [b] -> [(a,b)]
allPairs = undefined  -- YOUR CODE

foldActions :: Monad m => (a -> b -> m a) -> a -> [b] -> m a
foldActions = undefined  -- YOUR CODE
\end{codeanswer}
\end{minipage}\par

\task{4.5. Changed type question}
Infer \code{\xs -> bundle (map Just xs)} and describe its behavior. Explain
why its monad is more specific than the monad in \code{bundle}'s signature.
\answer{[Most general type and behavior.]}
